% from ega4-18-fr.tex
\oldpage[IV-4]{168}

Le raisonnement de (II, 7.3.8) prouve qu'il revient au même de dire que $f$ (supposé localement de type fini) est séparé et que, pour tout changement de base $Y''\to Y$, où $Y''$ est localement noethérien, l'image par $f_{(Y'')}:X_{(Y'')}\to Y''$ de toute partie fermée irréductible de $X''$ est fermée. Dire que $f$ est propre (pour $Y$ localement noethérien) signifie donc que $f$ est essentiellement propre et de type fini (II, 7.3.8); mais on rencontre des exemples importants de morphismes essentiellement propres qui ne sont pas de type fini (par exemple certains «préschémas de Picard» ou certains «préschémas de Néron-Severi» (chap. VI)). Il résulte évidemment de (18.10.19) que si $Y$ est localement noethérien, réduit et géométriquement unibranche, et si le morphisme $f:X\to Y$ est essentiellement propre et non ramifié, alors toute $Y$-section rationnelle de $f$ est partout définie.
\end{remark}

\begin{proposition}[18.10.21]
\label{IV.18.10.21}
Soient $f:X\to Y$ un morphisme localement de présentation finie, $x$ un point de $X$, $y=f(x)$ et supposons l'anneau $\sh{O}_{Y,y}$ intègre et géométriquement unibranche. Soit $r$ le nombre minimum de générateurs du $\sh{O}_{X,x}$-module $(\Omega^1_{X/Y})_x$ (égal au rang du $k(x)$-espace vectoriel $(\Omega^1_{X/Y})_x\otimes_{\sh{O}_{X,x}}k(x))$. Alors, pour que $f$ soit lisse au point $x$, il faut et il suffit que, si $y'$ est le point générique de l'unique composante irréductible de $Y$ contenant $y$, il existe une générisation $x'$ de $x$ telle que $f(x')=y'$ et $\dim_{x'}f^{-1}(y')\geq r$.
\end{proposition}

Si $f$ est lisse au point $x$, il l'est aussi dans un voisinage ouvert $U$ de $x$, donc en toute générisation $x'$ de $x$ et $\dim_{x'}(f^{-1}(f(x')))=r$ en chacun de ces points, en vertu de (17.10.2). Comme $f$ est en outre plat dans $U$, il existe une générisation de $x$ au-dessus de toute générisation de $y$ (2.3.4), ce qui prouve que la condition est nécessaire. Pour montrer qu'elle est suffisante, notons d'abord que par (17.5.1), il suffit de prouver que le morphisme déduit de $f$ par le changement de base $\Spec(\sh{O}_{Y,y})\to Y$ est lisse au point $x$; on peut donc supposer que $Y=\Spec(\sh{O}_{Y,y})$.

La question est locale sur $X$, donc ((16.4.22) et ($0_{\mathrm I}$, 5.2.2)) on peut supposer que $X$ est affine, et qu'il existe $r$ sections $t_i$ ($1\leq i\leq r$) de $\sh{O}_X$ au-dessus de $X$ telles que les sections $d_{X/Y}t_i$ engendrent $\Omega^1_{X/Y}$. Soit
\[
  g:X\to Z=Y[T_1,\ldots,T_r]=\mathbf V^r_Y
\]
le $Y$-morphisme correspondant à l'homomorphisme $\sh{O}_Y^r\to f_*(\sh{O}_X)$ de $\sh{O}_Y$-Modules défini par les $t_i$ (considérées comme sections de $f_*(\sh{O}_X)$ au-dessus de $Y$) (II, 1.2.7). On a la suite exacte (16.4.19)
\[
  g^*(\Omega^1_{Z/Y})\longrightarrow\Omega^1_{X/Y}\longrightarrow\Omega^1_{X/Z}\longrightarrow0.
\]
Comme les $dT_i$ ($1\leq i\leq r$) engendrent $\Omega^1_{Z/Y}$, il résulte de la définition de $g$ que l'homomorphisme $g^*(\Omega^1_{Z/Y})\to\Omega^1_{X/Y}$ est surjectif, donc $\Omega^1_{X/Z}=0$ et $g$ est par suite non ramifié (17.4.1). Nous allons voir que $g$ est étale au point $x$, et compte tenu de (17.3.8), cela prouvera que $f$ est lisse au point $x$. Comme $\sh{O}_{Y,y}$ est supposé intègre et géométriquement

